% Fitbauer: a Python-based open-source program for Mössbauer spectrum fitting
% Draft for submission to Hyperfine Interactions (Springer)
% Author: Jorge Sánchez Marcos
% Status: DRAFT — figures, validation section and author affiliation pending
%
% Compile with: pdflatex paper_hyperfine_interactions.tex
%
\documentclass[smallcondensed]{svjour3}
\smartqed
\usepackage[utf8]{inputenc}
\usepackage[T1]{fontenc}
\usepackage{amsmath}
\usepackage{amssymb}
\usepackage{graphicx}
\usepackage{hyperref}
\usepackage{booktabs}
\usepackage{doi}

\journalname{Interactions}  % Renamed from Hyperfine Interactions in 2024

\begin{document}

\title{Fitbauer: an open-source Python program for Mössbauer spectrum fitting
       with discrete components, hyperfine field distributions and magnetic
       relaxation models}

% TODO: complete with co-authors if applicable
\author{Jorge Sánchez Marcos}

\institute{%
  J. Sánchez Marcos \at
  Departamento de Química Física,
  Universidad Autónoma de Madrid,
  28049 Madrid, Spain \\
  \email{jorge.sanchezm@uam.es}
}

\date{Received: \today}

\maketitle

%----------------------------------------------------------------------
\begin{abstract}
We present Fitbauer, a free and open-source program written in Python for
the analysis of $^{57}$Fe Mössbauer spectra. The program combines a
graphical user interface built on PySide6~(Qt6) with a modular computational
core that is fully independent of the interface, enabling headless and
command-line use. Fitbauer provides the standard set of discrete components
(singlet, doublet, sextet) together with three magnetic relaxation models:
a phenomenological blocked-fraction model, the two-state dynamic model of
Blume and Tjon, and a Néel--Arrhenius model with lognormal particle-size
distribution for superparamagnetic systems. Quadrupole splitting is treated
either as a first-order rigid-pattern correction or through exact
diagonalisation of the combined magnetic-quadrupole Hamiltonian following
the approach of Kündig. Hyperfine-field distributions $P(B_\mathrm{HF})$ and
quadrupole-splitting distributions $P(\Delta E_\mathrm{Q})$ are obtained by
the Hesse--Rübartsch method with Tikhonov or Total Variation regularisation;
two-dimensional distributions $P(B_\mathrm{HF},\Delta E_\mathrm{Q})$,
$P(\mathrm{IS})$ and $P(\mathrm{IS},\Delta E_\mathrm{Q})$ are also
available. Statistical analysis includes weighted Poisson fitting,
reduced $\chi^2$, AIC, BIC, profile-likelihood asymmetric confidence
intervals, and bootstrap Monte Carlo errors. The program runs on Linux,
Windows and macOS under Python~3.11+ and is published under the Apache~2.0
licence at \url{https://github.com/sullymike/Mossbauer}.
\end{abstract}

\keywords{Mössbauer spectroscopy \and spectrum fitting software \and
hyperfine field distribution \and magnetic relaxation \and
Kündig quadrupole treatment \and Python}

%----------------------------------------------------------------------
\section{Introduction}
\label{sec:intro}

Mössbauer spectroscopy of $^{57}$Fe is a powerful technique for probing
local magnetic and electronic environments in solids. Reliable extraction
of hyperfine parameters requires fitting software that handles the full
variety of spectra encountered in practice: resolved sextets from ordered
phases, broad distributions arising from structural disorder or
superparamagnetism, relaxation lineshapes in nanoparticle systems, and
combinations thereof.

Several programs have been developed over the decades.
NORMOS~\cite{normos} is the long-standing commercial standard.
MossWinn~\cite{mosswinn2013} is a widely-used Windows program offering an
extensive database and advanced fitting methods.
SpectrRelax~\cite{spectrrelax2012} handles relaxation and distribution models
on Windows. Moessfit~\cite{moessfit2016} is a free, open-source Qt/C++
program focused on simultaneous multi-dataset fitting and the Maximum Entropy
Method. Vinda~\cite{vinda2016} provides Excel-based analysis.

Fitbauer is a new open-source Python implementation that fills a specific
niche: a cross-platform graphical program combining discrete components with
three distinct relaxation models, one-dimensional and two-dimensional
hyperfine-parameter distributions, exact Kündig quadrupole treatment, and a
complete statistical toolkit, all within a single installation requiring only
standard Python packages.

%----------------------------------------------------------------------
\section{Program overview}
\label{sec:overview}

\subsection{Architecture}
\label{sec:arch}

Fitbauer is organised in two independent layers. The computational core
(\texttt{core/}) consists of pure Python modules that depend only on NumPy
and SciPy; it has no dependency on any graphical toolkit. The graphical
interface (\texttt{gui/}) is built on PySide6~(Qt6) and Plotly, and calls
the core exclusively through well-defined API functions. This separation
allows the same physics engine to be used from the interactive GUI, from the
command-line interface (\texttt{mossbauer\_fit\_cli.py}), or from user
scripts.

\subsection{Supported data formats}
\label{sec:formats}

Two file formats are read natively: the modern WS5 format (XML-encoded
channel counts) and the legacy ADT format (plain integer count lists). When
a NORMOS sidecar \texttt{.RES} or \texttt{.PLT} file is present alongside
the spectrum, the folding point and Vmax are imported automatically as
initial values. Measurements and $\alpha$-Fe calibrations can also be
downloaded directly from a laboratory REST API from within the program.

\subsection{Typical workflow}
\label{sec:workflow}

After loading a spectrum the user inspects the folding point and velocity
axis (Vmax), selects a model (discrete or distribution), sets the initial
parameters, and launches the fit. The program displays the spectrum, model,
individual components and residual simultaneously. Results are saved as a
JSON session (complete, reproducible state), a plain-data file (velocity,
normalised counts, model, residual) and, optionally, a Markdown or PDF
report generated automatically.

%----------------------------------------------------------------------
\section{Physical models}
\label{sec:models}

\subsection{Folding and velocity axis}
\label{sec:folding}

The triangular Doppler waveform is folded at a fractional channel
(NORMOS convention). To avoid stretching the velocity scale when the two
edge channels are discarded, the full axis is constructed first on all $N$
channels and then trimmed at positions $[1\!:\!-1]$, consistent with the
treatment in NORMOS~\cite{normos}.

\subsection{Discrete spectral components}
\label{sec:discrete}

Up to six spectral components can be superimposed. Each component is one of
six types: singlet, doublet, sextet, and three relaxation types described in
Section~\ref{sec:relax}. For sextets the intensity pattern can be set to the
random-powder ratio $3:2:1$, left free ($I_1$, $I_2$, $I_3$ independent),
or parameterised by a single texture angle $t = \sin^2\!\theta$ following
the formula $3 : 4t/(2-t) : 1$.

\subsubsection{Line profile}
\label{sec:profile}

Two line profiles are available. The default Lorentzian profile is
\begin{equation}
  L(v;\,c,\gamma) = \frac{\gamma^2}{(v-c)^2 + \gamma^2},
\end{equation}
where $\gamma$ is the half-width at half-maximum (HWHM). The Voigt profile
convolves a Lorentzian of width $\gamma$ with a Gaussian of width $\sigma$,
evaluated via the Faddeeva function implemented in \texttt{scipy.special.wofz}.
The Voigt is normalised to its peak analytically, independently of the
velocity sampling.

\subsubsection{Quadrupole treatment}
\label{sec:quadrupole}

The conventional first-order treatment applies $\Delta E_\mathrm{Q}$ as a
rigid additive correction ($\pm\Delta E_\mathrm{Q}/2$ on the outer lines,
$\mp\Delta E_\mathrm{Q}/2$ on the inner lines). For cases where $\Delta
E_\mathrm{Q}$ is large or comparable to the magnetic splitting, an exact
Kündig treatment~\cite{kuendig1967} is available. The excited-state
Hamiltonian
\begin{equation}
  \hat{H} = \omega_L \hat{I}_z
  + \frac{\Delta E_\mathrm{Q}}{6}
    \left(3\hat{I}_{z'}^2 - \hat{I}(\hat{I}+1)\right)
  \label{eq:kuendig}
\end{equation}
is diagonalised numerically for each set of parameters, where $\omega_L$ is
the magnetic Zeeman splitting of the excited state, $z$ is along the
hyperfine field and $z'$ is along the principal axis of the electric-field
gradient. The angle $\beta$ between these two axes is either a free
parameter (\emph{Kündig fixed}) or integrated over all orientations by
20-point Gauss--Legendre quadrature on $\cos\beta \in [-1,1]$
(\emph{Kündig powder}). The asymmetry parameter $\eta$ of the EFG is taken
as zero (axially symmetric).

\subsubsection{Thick-absorber correction}
\label{sec:thick}

For thick absorbers the transmission is modelled as~\cite{margulies1961}
\begin{equation}
  T(v) = b + sv - C\!\left(1 - e^{-A_\mathrm{total}(v)/C}\right),
\end{equation}
where $b$ and $s$ are the baseline and slope, $A_\mathrm{total}(v)$ is the
total thin-absorber absorption, and $C$ is the saturation scale. In the
limit $C\to\infty$ the linear (thin-absorber) model is recovered. The
parameter $C$ can be fixed or fitted.

\subsection{Magnetic relaxation models}
\label{sec:relax}

Three relaxation components are available as regular spectral components and
can be combined freely with each other or with discrete phases.

\subsubsection{Phenomenological model (Relajacion)}
\label{sec:relajacion}

The phenomenological model interpolates linearly between a blocked-sextet
lineshape and an apparent superparamagnetic doublet:
\begin{equation}
  A(v) = A_0\!\left[f_\mathrm{eff}\,S_\mathrm{sex}(v)
         + (1 - f_\mathrm{eff})\,D_\mathrm{dbl}(v)\right],
\end{equation}
where the effective blocked fraction $f_\mathrm{eff}$ depends on both the
blocked-fraction parameter $f_\mathrm{b}$ and the relaxation rate
$\log_{10}\nu$ (s$^{-1}$). Because these two parameters are strongly
correlated, practical use requires fixing one while fitting the other.

\subsubsection{Two-state dynamic model (BlumeTjon)}
\label{sec:blumetjon}

The two-state model of Blume and Tjon~\cite{blumetjon1968} is implemented
for the $+B_\mathrm{HF} \leftrightarrow -B_\mathrm{HF}$ exchange process.
The spectrum is computed from the stochastic Liouville equation; the single
parameter $\log_{10}\nu$ drives the transition from slow (blocked sextet) to
fast (averaged doublet/singlet) relaxation through the intermediate
coalescence regime.

\subsubsection{Néel--Arrhenius model with particle-size distribution (NeelSize)}
\label{sec:neelsize}

The NeelSize component models a polydisperse nanoparticle assembly through a
lognormal diameter distribution $p(d;\,d_{50},\sigma_\ell)$. Each diameter
bin is converted to a volume $V = \pi d^3/6$ and a relaxation rate via the
Néel--Arrhenius equation~\cite{neel1949,brown1963}:
\begin{equation}
  \nu(d,T) = \frac{1}{\tau_0}
             \exp\!\left(-\frac{K_\mathrm{eff}\,V(d)}{k_\mathrm{B}\,T}\right).
\end{equation}
The total spectrum is the weighted sum of BlumeTjon spectra over the size
bins. The parameters are $T$, $\log_{10}K_\mathrm{eff}$, the median diameter
$d_{50}$, the lognormal width $\sigma_\ell$, $\log_{10}\tau_0$ and the
number of bins. A multi-temperature global-fit backend shares the physical
parameters across spectra while keeping the baseline, slope and depth local
to each spectrum.

%----------------------------------------------------------------------
\section{Hyperfine-parameter distributions}
\label{sec:distributions}

\subsection{\texorpdfstring{$P(B_\mathrm{HF})$ and $P(\Delta E_\mathrm{Q})$}{P(BHF) and P(dEQ)}}
\label{sec:pbhf}

One-dimensional distributions are obtained by the Hesse--Rübartsch
method~\cite{hesse1974}. A uniform grid of $N_\mathrm{bin}$ values is
defined between $B_\mathrm{min}$ and $B_\mathrm{max}$ (or
$\Delta E_{\mathrm{Q},\mathrm{min}}$ to $\Delta E_{\mathrm{Q},\mathrm{max}}$).
For each grid point the corresponding subspectrum is computed, forming the
response matrix $\mathbf{A}$. The distribution weights $\mathbf{p} \geq 0$
that minimise
\begin{equation}
  \chi^2 =
  \|\mathbf{y} - \mathbf{A}\mathbf{p}\|_w^2
  + \alpha\,\|\mathbf{D}^2\mathbf{p}\|^2
  \label{eq:tikhonov}
\end{equation}
are found by non-negative least squares, where $\|\cdot\|_w$ denotes the
Poisson-weighted norm and $\mathbf{D}^2$ is the second-difference matrix.
Alternatively, Total Variation (TV) regularisation, which better preserves
sharp edges in $P$, can be selected in the distribution panel.

Discrete ``sharp'' components (sextets, doublets, singleted) can be
superimposed on the distribution in the same fit; their columns are added to
$\mathbf{A}$ without regularisation.

\subsection{Two-dimensional distributions}
\label{sec:2d}

Two-dimensional distributions $P(B_\mathrm{HF}, \Delta E_\mathrm{Q})$,
$P(\mathrm{IS})$, $P(\mathrm{IS}, \Delta E_\mathrm{Q})$ and
$P(B_\mathrm{HF}, \mathrm{IS})$ are available. The 2D model minimises the
residual with independent smoothness penalties in each direction:
\begin{equation}
  \|\mathbf{y} - \mathbf{X}\mathbf{p}\|_w^2
  + \alpha_B\|\mathbf{D}^2_B\mathbf{p}\|^2
  + \alpha_Q\|\mathbf{D}^2_Q\mathbf{p}\|^2,
  \quad \mathbf{p} \geq 0.
\end{equation}
For a $30\times30$ grid the problem has 900 unknowns for typically
200--500 spectral channels and is therefore strongly underdetermined.
Users are cautioned to verify stability with respect to both regularisation
parameters and grid resolution before physical interpretation.

\subsection{Selection of the regularisation parameter}
\label{sec:lcurve}

The L-curve tool scans a wide range of $\alpha$ values and displays the
log--log plot of residual norm against roughness. Three suggestions are
provided: the maximum-curvature corner of the L-curve, a compromise value,
and the Generalised Cross-Validation (GCV) minimum~\cite{hansen1992}. A 1$\sigma$
uncertainty band around $P$ is drawn using the linearised covariance of the
weights (truncated at $P \geq 0$), and the reduced $\chi^2$ uses effective
degrees of freedom (trace of the hat matrix) rather than the raw number of
bins.

%----------------------------------------------------------------------
\section{Statistical analysis}
\label{sec:statistics}

\subsection{Fitting algorithm}
\label{sec:fitalgo}

The discrete fit uses the Trust-Region Reflective (TRF) algorithm
(\texttt{scipy.optimize.least\_squares}) with a deterministic
nine-start multi-start strategy to reduce dependence on the initial
parameters. An optional prior pass with Differential Evolution
(\texttt{scipy.optimize.differential\_evolution}) can be enabled for
complex multi-sextet models. Robust loss functions (linear,
soft-$L_1$, Huber) are available to reduce the influence of outlier
channels.

Two likelihood modes are provided: Gaussian ($\chi^2$, default) and
Poisson IRLS, where the weights are iteratively rescaled from the model
prediction rather than from the data.

\subsection{Parameter uncertainties}
\label{sec:errors}

Standard $1\sigma$ uncertainties are estimated from the Jacobian-based
covariance matrix. For strongly non-linear problems or parameters near a
physical boundary, asymmetric confidence intervals are computed by profile
likelihood: for each free parameter the objective is profiled over a dense
grid, the complementary parameters are re-optimised at each grid point, and
the $\Delta\chi^2 = 1$ (1$\sigma$) and $\Delta\chi^2 = 4$ (2$\sigma$)
crossing points are located by interpolation. Bootstrap Monte Carlo errors
are also available for discrete models: synthetic spectra are generated by
adding Poisson noise to the current model and refitted; the standard
deviation over replicas is reported.

\subsection{Model selection}
\label{sec:modelsel}

After fitting, the program reports the reduced $\chi^2$, Akaike
Information Criterion (AIC) and Bayesian Information Criterion (BIC).
Linear constraints between parameters (e.g.\ equal line widths, fixed
intensity ratios) can be specified by the user; physical presets for
common relationships (random-powder $3:2:1$ intensities, equal widths
across components) are available as one-click shortcuts.

Three residual diagnostics are shown automatically: the lag-1
autocorrelation, the runs test statistic (sensitive to systematic trends),
and an antisymmetric correlation coefficient (sensitive to an incorrect
folding point).

%----------------------------------------------------------------------
\section{Implementation}
\label{sec:implementation}

Fitbauer is written in Python~3.11+ and requires NumPy, SciPy, Matplotlib,
Plotly and PySide6. No compiled extension modules are needed; the program
installs as a plain Python package (\texttt{pip install -r requirements.txt}).
Pre-built ZIP archives for Linux, Windows and macOS are distributed via
GitHub Releases. The graphical interface runs in three languages (Spanish,
English, French); the help system with ~30 thematic chapters is available in
all three. An automated update mechanism checks GitHub Releases on startup.

The full test suite (\texttt{pytest}) covers the physics core, the fit
engine, folding, the headless session API and the Qt GUI (via
\texttt{pytest-qt} with \texttt{QT\_QPA\_PLATFORM=offscreen}). Continuous
integration runs on Python~3.12 / PySide6 via GitHub Actions.

%----------------------------------------------------------------------
\section{Validation}
\label{sec:validation}

% TODO: insert one or two real-data examples
% Suggested structure:
%   7.1  alpha-Fe calibration spectrum (sextet, verify LINE_POS_33T)
%   7.2  Magnetite Fe3O4 (two sextets A+B, illustrate workflow)
%   7.3  Disordered iron oxide or nanoparticle sample (P(BHF), L-curve)
%   7.4  Optional: nanoparticles with NeelSize (temperature series)
%
% Each example should include: spectrum + fit figure, table of parameters
% compared to reference values from the literature.
%
\textbf{[PENDING — figures and validation data to be added by the author]}

%----------------------------------------------------------------------
\section{Conclusions}
\label{sec:conclusions}

Fitbauer is a fully open-source Mössbauer fitting program that combines,
within a single cross-platform Qt application, standard discrete components,
three magnetic relaxation models (phenomenological, Blume--Tjon two-state,
Néel--Arrhenius with particle-size distribution), exact Kündig quadrupole
diagonalisation, regularised hyperfine-parameter distributions in one and
two dimensions, and a complete statistical toolkit including profile-likelihood
confidence intervals and bootstrap Monte Carlo errors. The program targets
users who require the full range from simple calibration fits to advanced
distribution and relaxation analysis, without resorting to multiple
specialised tools.

\subsection*{Availability}

Fitbauer version~4.7.1 is freely available under the Apache~2.0 licence
at \url{https://github.com/sullymike/Mossbauer}. Pre-built archives
for Linux, Windows and macOS are distributed via GitHub Releases.

\subsection*{Acknowledgements}

% TODO: add funding acknowledgements if applicable.

%----------------------------------------------------------------------
% REFERENCES
% Hyperfine Interactions uses numbered Vancouver-style references.
%----------------------------------------------------------------------
\begin{thebibliography}{99}

\bibitem{normos}
  Brand, R.A.: NORMOS Mössbauer fitting program.
  Universität Duisburg (1987); WissEl GmbH, Düsseldorf (current distributor).

\bibitem{mosswinn2013}
  Klencsár, Z.: MossWinn — methodological advances in the field of
  Mössbauer data analysis.
  Hyperfine Interact. \textbf{217}, 117--126 (2013).
  \doi{10.1007/s10751-012-0732-2}

\bibitem{spectrrelax2012}
  Matsnev, M.E., Rusakov, V.S.: SpectrRelax: an application for Mössbauer
  spectra modeling and fitting.
  AIP Conf. Proc. \textbf{1489}, 178--185 (2012).
  \doi{10.1063/1.4759488}

\bibitem{moessfit2016}
  Kamusella, S., Klauss, H.-H.: Moessfit: a free Mössbauer fitting program.
  Hyperfine Interact. \textbf{237}, 82 (2016).
  \doi{10.1007/s10751-016-1247-z}

\bibitem{vinda2016}
  Gunnlaugsson, H.P.: Spreadsheet based analysis of Mössbauer spectra.
  Hyperfine Interact. \textbf{237}, 79 (2016).
  \doi{10.1007/s10751-016-1271-z}

\bibitem{hesse1974}
  Hesse, J., Rübartsch, A.: Model independent evaluation of overlapped
  Mössbauer spectra.
  J. Phys. E: Sci. Instrum. \textbf{7}, 526--532 (1974).
  \doi{10.1088/0022-3735/7/7/015}

\bibitem{kuendig1967}
  Kündig, W.: Evaluation of the electric-field gradient and the asymmetry
  parameter $\eta$ in Mössbauer spectra.
  Nucl. Instrum. Methods \textbf{48}, 219--228 (1967).
  \doi{10.1016/0029-554X(67)90513-1}

\bibitem{margulies1961}
  Margulies, S., Ehrman, J.R.: Transmission and line broadening of resonance
  radiation incident on a resonance absorber.
  Nucl. Instrum. Methods \textbf{12}, 131--137 (1961).
  \doi{10.1016/0029-554X(61)90210-X}

\bibitem{blumetjon1968}
  Blume, M., Tjon, J.A.: Mössbauer spectra in a fluctuating environment.
  Phys. Rev. \textbf{165}, 446--456 (1968).
  \doi{10.1103/PhysRev.165.446}

\bibitem{neel1949}
  Néel, L.: Théorie du traînage magnétique des ferromagnétiques en grains
  fins avec applications aux terres cuites.
  Ann. Géophys. \textbf{5}, 99--136 (1949).

\bibitem{brown1963}
  Brown, W.F.: Thermal fluctuations of a single-domain particle.
  Phys. Rev. \textbf{130}, 1677--1686 (1963).
  \doi{10.1103/PhysRev.130.1677}

\bibitem{hansen1992}
  Hansen, P.C.: Analysis of discrete ill-posed problems by means of the
  L-curve.
  SIAM Rev. \textbf{34}, 561--580 (1992).
  \doi{10.1137/1034115}

\bibitem{rancourt1991}
  Rancourt, D.G., Ping, J.Y.: Voigt-based methods for arbitrary-shape static
  hyperfine parameter distributions in Mössbauer spectroscopy.
  Nucl. Instrum. Methods B \textbf{58}, 85--97 (1991).
  \doi{10.1016/0168-583X(91)95681-3}

\end{thebibliography}

\end{document}
